\selectlanguage{ngerman}
\frenchspacing
\begin{ledgroupsized}[r]{114mm}
\footnotesize
\pstart
\noindent
\textbf{Überlieferung:}
\pend
\end{ledgroupsized}
\begin{ledgroupsized}[r]{114mm}
\footnotesize
\pstart
\parindent -6mm
\makebox[6mm][l]{\textit{L}}%
Notiz:
LH XXXV~6,~23 Bl.~3.
Ein Blatt 8\textsuperscript{o} (italienisches Papier);
Ränder zum Teil beschnitten.
Eindrittel Seite im unteren Bereich von Bl.~\href{\Scan{LH_35_06_23_003r}}{3~r\textsuperscript{o}}.
Bl.~\href{\Scan{LH_35_06_23_003r}}{3~r\textsuperscript{o}} und \href{\Scan{LH_35_06_23_003v}}{3~v\textsuperscript{o}} überliefern außerdem die Aufzeichnung N.\;\ref{58656-A}.
\pend
\end{ledgroupsized}
\vspace{5mm}
\begin{ledgroup}
\footnotesize
\pstart
\noindent
\textbf{Datierungsgründe:}
Die gemeinsame Überlieferung der Notiz N.\;\ref{58656-B} und der Aufzeichnung N.\;\ref{58656-A} rechtfertigt eine gemeinsame Datierung beider Texte. 
Letztere Aufzeichnung N.\;\ref{58656-A} ist aller Wahrscheinlichkeit nach im Sommer bis Herbst 1689 entstanden, im Zusammenhang mit Leibnizens Entwürfe zur \cite{01354}\textit{Dynamica de potentia} (siehe die Datierung auf Grund inhaltlicher Betrachtungen auf S.~\pageref{58656-A_dat}).
Mangels weiterer, inhaltlicher Anhaltspunkte für die genauere Datierung der vorliegenden Notiz N.\;\ref{58656-B} wird obengenannte Datierung übernommen.
\pend
\end{ledgroup}
\selectlanguage{latin}
\vspace{8mm}
\pstart
\noindent
\href{\Scan{LH_35_06_23_003r}}%
{\lbrack3~r\textsuperscript{o}\rbrack}
Si in figuram \textit{ABC} ducamus figuram \textit{ABDEF}, eadem est 
\edtext{difficultas in via\indexS{via centri gravitatis}}{%
\lemma{difficultas in}%
\Bfootnote{%
\textit{(1)}~figura
\textit{(2)}~via~\textit{L}}}
centri gravitatis\indexS{centrum gravitatis} \edlabel{35_06_23_003r_1}ductus
\edtext{invenienda, et}{%
\lemma{invenienda,}%
\Bfootnote{%
\textit{(1)}~neque
\textit{(2)}~et~\textit{L}}}
licet ponamus \textit{ABC} et \textit{FABDEF} esse utrinque congruas quamlibet suis partibus,
hinc quidem sequitur viam\indexS{via centri gravitatis} esse in linea recta\lbrack,\rbrack\ sed ultimum punctum
\edtext{postulat centrum}{%
\lemma{postulat}%
\Bfootnote{%
\textit{(1)}~viam
\textit{(2)}~centrum~\textit{L}}}
gravitatis omnium rectangulorum ut \textit{EG}, ubi relicta sunt, non vero centrum gravitatis\indexS{centrum gravitatis}
superficiei solidi generati.\indexS{centrum gravitatis superficiei solidi generati}
\pend

\begin{edfigure}%
\hfill
\includegraphics[scale=1]{%
gesamttex/edit_VIII-4/images/LH_35_06_23_003r_d1.pdf}
\hfill%
\includegraphics[scale=1]{%
gesamttex/edit_VIII-4/images/LH_35_06_23_003r_d2.pdf}
\hfill
\includegraphics[scale=1]{%
gesamttex/edit_VIII-4/images/LH_35_06_23_003r_d3.pdf}
\hspace{1cm}
\edcaption{\lbrack\textit{Fig.~1}\rbrack 
\hspace*{30mm}%
\lbrack\textit{Fig.~2}\rbrack
\hspace*{35mm}%
\edtext{\lbrack\textit{Fig.~3}\rbrack}{%
\lemma{\lbrack\textit{Fig.~3}\rbrack}%
\Cfootnote[nonum]{Zwei gestrichene Entwürfe zum Diagramm, davon einer auf Höhe des Textes von \refpassage{35_06_23_003r_1}{35_06_23_003r_1}, werden nicht wiedergegeben.}}}
\end{edfigure}
